% GENERATED FILE. Edit canonical lessons, then run npm run generate:books.
% canonical-source-hash: fnv1a64:74410e018657a3e0
% canonical-lessons: SA-C309-vrttantah, SA-C309-prasarayati, SA-C309-akarnayati, SA-C309-mukhyah, SA-C309-vartaharah, SA-R309-review-302-309

\chapter{The Morning News}
\label{ch:sa-the-morning-news}
\hlchaptermodality{309}

\begin{chapteropening}
\textbf{By the end of this chapter:} \emph{I can say what the news reported.}

Complete the last lesson of chapter 309: I can say what the news reported.
\end{chapteropening}

\section[\texorpdfstring{\sk{वृत्तान्तः} (vṛttāntaḥ)}{वृत्तान्तः (vṛttāntaḥ)}]{\sk{वृत्तान्तः} (vṛttāntaḥ) — an account; a report}

\label{lesson:SA-C309-vrttantah}



\begin{quote}
Before the new one: say the Sanskrit for a donor, then the Sanskrit for a helper.
\end{quote}

\subsection*{You'll want to know: \sk{वृत्तान्तः}}
\textbf{\sk{वृत्तान्तः}} — \emph{vṛttāntaḥ} — \textquotedblleft{}an account; a report\textquotedblright{}.

\begin{quote}
\sk{पत्रिकायां} \sk{जलौघस्य} \sk{वृत्तान्तः} \sk{अस्ति।}
\end{quote}

\begin{quote}
\emph{patrikāyāṃ jalaughasya vṛttāntaḥ asti.}
\end{quote}

\begin{quote}
There is an account of the flood in the newspaper.
\end{quote}

A masculine noun.

\begin{cousinweb}
\textbf{\sk{वृत्तान्तः}} (\emph{vṛttāntaḥ}) is built from \emph{vṛtta}, \textquotedblleft{}what has happened\textquotedblright{}, and \emph{anta}, \textquotedblleft{}an end\textquotedblright{}; MW gives \emph{vṛttānta} as an occurrence or event and also as tidings, a report, a tale (MW).
\end{cousinweb}

\begin{grammarlens}[title={What you've built}]
One more word for telling what happened.
\end{grammarlens}

\subsection*{Your turn}
\begin{itemize}
\raggedright
  \item \emph{Say it:} \emph{vṛttāntaḥ}
  \item \emph{Say it:} \emph{vṛttāntaḥ}, once more
  \item \emph{Say it:} \emph{patrikāyāṃ jalaughasya vṛttāntaḥ asti.} — There is an account of the flood in the newspaper.
  \item \emph{Recall:} say the Sanskrit for a donor, then the Sanskrit for a helper, then say \emph{vṛttāntaḥ} again
\end{itemize}

\begin{tcolorbox}[breakable,colback=teal!4,colframe=teal!35!black,title={Before you move on}]
What does \sk{वृत्तान्तः} mean? (\textquotedblleft{}An account; a report\textquotedblright{}.) Say it once more.
\end{tcolorbox}
\section[\texorpdfstring{\sk{प्रसारयति} (prasārayati)}{प्रसारयति (prasārayati)}]{\sk{प्रसारयति} (prasārayati) — broadcasts; spreads out}

\label{lesson:SA-C309-prasarayati}



\begin{quote}
Before the new one: say the Sanskrit for a helper, then the Sanskrit for an account.
\end{quote}

\subsection*{You'll want to know: \sk{प्रसारयति}}
\textbf{\sk{प्रसारयति}} — \emph{prasārayati} — \textquotedblleft{}broadcasts; spreads out\textquotedblright{}.

\begin{quote}
\sk{आकाशवाणी} \sk{वार्तां} \sk{प्रसारयति।}
\end{quote}

\begin{quote}
\emph{ākāśavāṇī vārtāṃ prasārayati.}
\end{quote}

\begin{quote}
The radio broadcasts the news.
\end{quote}

Past (it …): \sk{आकाशवाणी} \sk{वार्तां} \sk{प्रसारितवती}, \emph{ākāśavāṇī vārtāṃ prasāritavatī}.

\begin{cousinweb}
\textbf{\sk{प्रसारयति}} (\emph{prasārayati}) is built from \emph{pra} and the causative of the root √sṛ, \textquotedblleft{}to run, flow\textquotedblright{}; MW gives \emph{pra-sārayati} as to stretch out, extend, spread out, and to broadcast is its modern use (MW).
\end{cousinweb}

\begin{grammarlens}[title={What you've built}]
One more word for telling what happened.
\end{grammarlens}

\subsection*{Your turn}
\begin{itemize}
\raggedright
  \item \emph{Say it:} \emph{prasārayati}
  \item \emph{Say it:} \emph{prasārayati}, once more
  \item \emph{Say it:} \emph{ākāśavāṇī vārtāṃ prasārayati.} — The radio broadcasts the news.
  \item \emph{Recall:} say the Sanskrit for a helper, then the Sanskrit for an account, then say \emph{prasārayati} again
\end{itemize}

\begin{tcolorbox}[breakable,colback=teal!4,colframe=teal!35!black,title={Before you move on}]
What does \sk{प्रसारयति} mean? (\textquotedblleft{}Broadcasts; spreads out\textquotedblright{}.) Say it once more.
\end{tcolorbox}
\section[\texorpdfstring{\sk{आकर्णयति} (ākarṇayati)}{आकर्णयति (ākarṇayati)}]{\sk{आकर्णयति} (ākarṇayati) — listens to}

\label{lesson:SA-C309-akarnayati}



\begin{quote}
Before the new one: say the Sanskrit for an account, then the Sanskrit for broadcasts.
\end{quote}

\subsection*{You'll want to know: \sk{आकर्णयति}}
\textbf{\sk{आकर्णयति}} — \emph{ākarṇayati} — \textquotedblleft{}listens to\textquotedblright{}.

\begin{quote}
\sk{पिता} \sk{प्रातः} \sk{वार्ताम्} \sk{आकर्णयति।}
\end{quote}

\begin{quote}
\emph{pitā prātaḥ vārtām ākarṇayati.}
\end{quote}

\begin{quote}
Father listens to the news in the morning.
\end{quote}

Past (he …): \sk{सः} \sk{वार्ताम्} \sk{आकर्णितवान्}, \emph{saḥ vārtām ākarṇitavān}.

\begin{cousinweb}
\textbf{\sk{आकर्णयति}} (\emph{ākarṇayati}) is built on \emph{ā} and \emph{karṇa}, \textquotedblleft{}the ear\textquotedblright{}: MW gives \emph{ākarṇayati} as a verb made from them, to give ear to, listen (MW).
\end{cousinweb}

\begin{grammarlens}[title={What you've built}]
One more word for telling what happened.
\end{grammarlens}

\subsection*{Your turn}
\begin{itemize}
\raggedright
  \item \emph{Say it:} \emph{ākarṇayati}
  \item \emph{Say it:} \emph{ākarṇayati}, once more
  \item \emph{Say it:} \emph{pitā prātaḥ vārtām ākarṇayati.} — Father listens to the news in the morning.
  \item \emph{Recall:} say the Sanskrit for an account, then the Sanskrit for broadcasts, then say \emph{ākarṇayati} again
\end{itemize}

\begin{tcolorbox}[breakable,colback=teal!4,colframe=teal!35!black,title={Before you move on}]
What does \sk{आकर्णयति} mean? (\textquotedblleft{}Listens to\textquotedblright{}.) Say it once more.
\end{tcolorbox}
\section[\texorpdfstring{\sk{मुख्यः} (mukhyaḥ)}{मुख्यः (mukhyaḥ)}]{\sk{मुख्यः} (mukhyaḥ) — main; chief}

\label{lesson:SA-C309-mukhyah}



\begin{quote}
Before the new one: say the Sanskrit for broadcasts, then the Sanskrit for listens to.
\end{quote}

\subsection*{You'll want to know: \sk{मुख्यः}}
\textbf{\sk{मुख्यः}} — \emph{mukhyaḥ} — \textquotedblleft{}main; chief\textquotedblright{}.

\begin{quote}
\sk{अद्य} \sk{मुख्या} \sk{वार्ता} \sk{जलौघः} \sk{अस्ति।}
\end{quote}

\begin{quote}
\emph{adya mukhyā vārtā jalaughaḥ asti.}
\end{quote}

\begin{quote}
Today the main news is the flood.
\end{quote}

\begin{cousinweb}
\textbf{\sk{मुख्यः}} (\emph{mukhyaḥ}) is built on \emph{mukha}, \textquotedblleft{}mouth, face, front\textquotedblright{}; MW gives \emph{mukhya} as being at the head, first, principal, chief (MW).
\end{cousinweb}

\begin{grammarlens}[title={What you've built}]
One more word for telling what happened.
\end{grammarlens}

\subsection*{Your turn}
\begin{itemize}
\raggedright
  \item \emph{Say it:} \emph{mukhyaḥ}
  \item \emph{Say it:} \emph{mukhyaḥ}, once more
  \item \emph{Say it:} \emph{adya mukhyā vārtā jalaughaḥ asti.} — Today the main news is the flood.
  \item \emph{Recall:} say the Sanskrit for broadcasts, then the Sanskrit for listens to, then say \emph{mukhyaḥ} again
\end{itemize}

\begin{tcolorbox}[breakable,colback=teal!4,colframe=teal!35!black,title={Before you move on}]
What does \sk{मुख्यः} mean? (\textquotedblleft{}Main; chief\textquotedblright{}.) Say it once more.
\end{tcolorbox}
\section[\texorpdfstring{\sk{वार्ताहरः} (vārtāharaḥ)}{वार्ताहरः (vārtāharaḥ)}]{\sk{वार्ताहरः} (vārtāharaḥ) — a news reporter}

\label{lesson:SA-C309-vartaharah}



\begin{quote}
Before the new one: say the Sanskrit for listens to, then the Sanskrit for main.
\end{quote}

\subsection*{You'll want to know: \sk{वार्ताहरः}}
\textbf{\sk{वार्ताहरः}} — \emph{vārtāharaḥ} — \textquotedblleft{}a news reporter\textquotedblright{}.

\begin{quote}
\sk{वार्ताहरः} \sk{यामिकं} \sk{पृच्छति।}
\end{quote}

\begin{quote}
\emph{vārtāharaḥ yāmikaṃ pṛcchati.}
\end{quote}

\begin{quote}
The reporter questions the watchman.
\end{quote}

A masculine noun.

\begin{cousinweb}
\textbf{\sk{वार्ताहरः}} (\emph{vārtāharaḥ}) is built from \emph{vārtā}, \textquotedblleft{}news\textquotedblright{} (MW writes \emph{vārttā}), and \emph{hara}, \textquotedblleft{}carrying\textquotedblright{}, from the root √hṛ; MW gives it as a carrier of tidings, a messenger, and a news reporter is its modern use (MW).
\end{cousinweb}

\begin{grammarlens}[title={What you've built}]
One more word for telling what happened.
\end{grammarlens}

\subsection*{Your turn}
\begin{itemize}
\raggedright
  \item \emph{Say it:} \emph{vārtāharaḥ}
  \item \emph{Say it:} \emph{vārtāharaḥ}, once more
  \item \emph{Say it:} \emph{vārtāharaḥ yāmikaṃ pṛcchati.} — The reporter questions the watchman.
  \item \emph{Recall:} say the Sanskrit for listens to, then the Sanskrit for main, then say \emph{vārtāharaḥ} again
\end{itemize}

\begin{tcolorbox}[breakable,colback=teal!4,colframe=teal!35!black,title={Before you move on}]
What does \sk{वार्ताहरः} mean? (\textquotedblleft{}A news reporter\textquotedblright{}.) Say it once more.
\end{tcolorbox}
\section[(dialogue)]{Looking back: from The Water Rose to The Morning News}

\label{lesson:SA-R309-review-302-309}



\begin{quote}
Say the Sanskrit for a flood and floods.
\end{quote}

\subsection*{Your turn}
Say these aloud:

\begin{itemize}
\raggedright
  \item a flood, floods, a current, goes under, by degrees
  \item a famine, an epidemic, wasted away, spreads through, withers
  \item a heap, in distress, shouts, laments, brings out
  \item rescues, life, daring, brings down, saves
  \item a camp, shelter, a donation, a refugee, comforts
  \item a separation, a tear, meets, orphaned, safety
  \item a group, does a kindness, shares out, a donor, a helper
  \item an account, broadcasts, listens to, main, a news reporter
\end{itemize}

\begin{tcolorbox}[breakable,colback=teal!4,colframe=teal!35!black,title={Before you move on}]
A flood? (\textbf{\sk{जलौघः}}, \emph{jalaughaḥ}.) Floods? (\textbf{\sk{प्लावयति}}, \emph{plāvayati}.) A current? (\textbf{\sk{प्रवाहः}}, \emph{pravāhaḥ}.) Goes under? (\textbf{\sk{निमज्जति}}, \emph{nimajjati}.) By degrees? (\textbf{\sk{क्रमशः}}, \emph{kramaśaḥ}.) A famine? (\textbf{\sk{दुर्भिक्षम्}}, \emph{durbhikṣam}.) An epidemic? (\textbf{\sk{महामारी}}, \emph{mahāmārī}.) Wasted away? (\textbf{\sk{क्षीणः}}, \emph{kṣīṇaḥ}.) Spreads through? (\textbf{\sk{व्याप्नोति}}, \emph{vyāpnoti}.) Withers? (\textbf{\sk{म्लायति}}, \emph{mlāyati}.) A heap? (\textbf{\sk{राशिः}}, \emph{rāśiḥ}.) In distress? (\textbf{\sk{आर्तः}}, \emph{ārtaḥ}.) Shouts? (\textbf{\sk{क्रोशति}}, \emph{krośati}.) Laments? (\textbf{\sk{विलपति}}, \emph{vilapati}.) Brings out? (\textbf{\sk{निःसारयति}}, \emph{niḥsārayati}.) Rescues? (\textbf{\sk{उद्धरति}}, \emph{uddharati}.) Life? (\textbf{\sk{प्राणाः}}, \emph{prāṇāḥ}.) Daring? (\textbf{\sk{साहसम्}}, \emph{sāhasam}.) Brings down? (\textbf{\sk{अवतारयति}}, \emph{avatārayati}.) Saves? (\textbf{\sk{त्रायते}}, \emph{trāyate}.) A camp? (\textbf{\sk{शिबिरम्}}, \emph{śibiram}.) Shelter? (\textbf{\sk{आश्रयः}}, \emph{āśrayaḥ}.) A donation? (\textbf{\sk{दानम्}}, \emph{dānam}.) A refugee? (\textbf{\sk{शरणार्थी}}, \emph{śaraṇārthī}.) Comforts? (\textbf{\sk{सान्त्वयति}}, \emph{sāntvayati}.) A separation? (\textbf{\sk{वियोगः}}, \emph{viyogaḥ}.) A tear? (\textbf{\sk{अश्रु}}, \emph{aśru}.) Meets? (\textbf{\sk{मिलति}}, \emph{milati}.) Orphaned? (\textbf{\sk{अनाथः}}, \emph{anāthaḥ}.) Safety? (\textbf{\sk{क्षेमः}}, \emph{kṣemaḥ}.) A group? (\textbf{\sk{समूहः}}, \emph{samūhaḥ}.) Does a kindness? (\textbf{\sk{उपकरोति}}, \emph{upakaroti}.) Shares out? (\textbf{\sk{विभजति}}, \emph{vibhajati}.) A donor? (\textbf{\sk{दाता}}, \emph{dātā}.) A helper? (\textbf{\sk{सहायः}}, \emph{sahāyaḥ}.) An account? (\textbf{\sk{वृत्तान्तः}}, \emph{vṛttāntaḥ}.) Broadcasts? (\textbf{\sk{प्रसारयति}}, \emph{prasārayati}.) Listens to? (\textbf{\sk{आकर्णयति}}, \emph{ākarṇayati}.) Main? (\textbf{\sk{मुख्यः}}, \emph{mukhyaḥ}.) A news reporter? (\textbf{\sk{वार्ताहरः}}, \emph{vārtāharaḥ}.)
\end{tcolorbox}
